\section{Threats to Validity}\label{sec:threats}

\paragraph*{Construct validity}
We measure dependence by perturbing each operation at \numInfluenceInputs{} inputs, which can miss a
dependence that shows only at other inputs. Rerun with \numGcInputs{} inputs and \numGcPerturbations{}
perturbations, the procedure finds \numGcDepsNew{} dependences
beyond the \numGcDepsBefore{} already measured (\numGcDepsNewPct{}), all inside QDG slices, and no rule's
mean recall moves by more than \numGcRecallShiftQ{}. Missed dependences cannot affect soundness, which is judged by the oracle's
fidelity at a threshold of $1 - 10^{-6}$, far above its error of \numOracleErr{}, and the QDG's soundness
for all inputs is proved. The other rules are judged
at $\ket{0\dots0}$ and on \numRobDraws{} draws per input regime, so they could be wrong at other inputs.
Containing the fault site is a proxy for helping a developer, and we ran no study with developers. The
profile predicates of the LintQ comparison are ours, fixed before the comparison, and show that a profile
exposes a confirmed bug, not that a developer would see it.

\paragraph*{Internal validity}
We recorded every choice of our primary reading of Yu and Zhao's prose before the runs. The reading in which a gate only reads its control, added after the results, is unsound on \numYzRoUnsound{} pairs, so claims
about the QPDG hold for the reading named. Rounding at the $10^{-9}$ thresholds of state-based pruning
flips some decisions on one Bugs4Q program, while \numQRepeatRuns{} fresh runs reproduce all
\numQRepeatPairs{} slices of corpus~A.

\paragraph*{External validity}
Exact simulation limits corpus~A to \numMaxQubits{} qubits, corpus~L reaches
\numCorpusLMaxQubits{} qubits, and Theorem~\ref{thm:sound} covers any width. The circuits come from
benchmarks, and the real bugs are few and small. The rules are stated on circuits and carry over to
other languages, except that a repeat-until-success loop in Q\#~\cite{svore2018qsharp} has no single
circuit and a static slicer must over-approximate its iterations. Seeded faults may differ from
real ones, although \numMutPublished{} of our \numMutOperators{} operators come from published tools and
fault models. Times are measured on one machine, and LintQ's and QChecker's are those their papers report on
theirs, so the comparison of times is indicative. \tool{}'s OpenQASM reader is written in Python and dominates its
time on the largest circuits.
